
\section{Related work}
\label{sec:related}

% Define colors for checkmark and cross
\definecolor{GreenCheck}{rgb}{0.2, 0.6, 0.2}
\definecolor{RedCross}{rgb}{0.8, 0.2, 0.2}

% \checkmark and \times MUST be in math mode ($...$)
\newcommand{\cmark}{\textcolor{GreenCheck}{$\checkmark$}}%
\newcommand{\xmark}{\textcolor{RedCross}{$\times$}}%



\begin{table}[!ht]
\centering
\caption{ \textbf{Comparison of benchmark features.}
A checkmark (\cmark) indicates that the benchmark explicitly includes the corresponding evaluation dimension or metric, while a cross (\xmark) indicates that it does not. 
In Core Dimensions, Perception(Percept); Prediction(Pred); Planning(Plan); Execution(Exec); Generalization(General); in Metrics Design, Video Quality(VQ); Instruction Understanding(IU); Physical Law(PL); Planning Reasoning(PR); Execution Accuracy(EA).
}

\label{tab:final_comparison}
\setlength{\tabcolsep}{2.2pt} % 控制列间距，可根据需要调整
\scriptsize
\begin{tabular}{@{}lcccccccccc@{}}
\toprule
& \multicolumn{5}{c}{{Core Dimensions}} &
  \multicolumn{5}{c}{{Metrics Design}} \\
\cmidrule(lr){2-6}\cmidrule(lr){7-11}
{Benchmark} &
{Percept} & {Pred} & {Plan} & {Exec} & {General} &
{VQ} & {IU} & {PL} & {PR} & {EA} \\
\midrule

Physics\textnormal{-}IQ~\cite{motamed2025generative}         & \cmark & \cmark & \xmark & \xmark & \xmark & \xmark & \xmark & \cmark & \xmark & \xmark \\
PhyGenBench~\cite{meng2024towards}                        & \cmark & \cmark & \xmark & \xmark & \xmark & \xmark & \xmark & \cmark & \xmark & \xmark \\

T2V\textnormal{-}CompBench~\cite{sun2025t2v}        & \cmark & \cmark & \xmark & \xmark & \xmark & \xmark & \cmark & \cmark & \xmark & \xmark \\
VBench\textnormal{-}2.0~\cite{zheng2025vbench}         & \cmark & \cmark & \xmark & \xmark & \cmark & \cmark & \cmark & \cmark & \xmark & \xmark \\

WorldModelBench~\cite{li2025worldmodelbench}                    & \cmark & \cmark & \xmark & \xmark & \cmark & \xmark & \cmark & \cmark & \xmark & \xmark \\
EWMBench~\cite{yue2025ewmbench}                           & \cmark & \cmark & \xmark & \xmark & \xmark & \xmark & \cmark & \cmark & \xmark & \xmark \\

{Ours}        & \cmark & \cmark & \cmark & \cmark & \cmark & \cmark & \cmark & \cmark & \cmark & \cmark \\

\bottomrule
\end{tabular}
\end{table}




The evaluation of generative world models and video generation models is rapidly evolving. Existing benchmarks can be grouped into: (1) those targeting general video quality, (2) those assessing physical reasoning, and (3) benchmarks for holistic world model evaluation. 

\textbf{Video Generation Benchmarks.} Early benchmarks mainly assessed visual fidelity and temporal coherence. TC-Bench~\cite{feng2024tc} studies temporal compositionality with transition-aware metrics. \cite{ling2025vmbench} measures perception-aligned motion quality. T2V-CompBench~\cite{sun2025t2v} evaluates compositionality via MLLM- and tracking-based methods, challenging metrics like FVD. \cite{liu2024evalcrafter,ji2024t2vbench,chi2024eva} introduce tasks related to controllability and anticipation. VBench-2.0~\cite{zheng2025vbench} focuses on “intrinsic faithfulness,” adding \textit{Physics} and \textit{Commonsense} dimensions using VLM/LLM evaluators.

\textbf{Physical Reasoning Benchmarks.} Physics-IQ~\cite{motamed2025generative} introduces physics-based diagnostics. PhyGenBench~\cite{meng2024towards} targets semantic and physical plausibility. PhysBench~\cite{chow2025physbench} evaluates VLM physical reasoning via QA tasks rather than pixel metrics. Videophy~\cite{bansal2024videophy} benchmarks physical commonsense using rule-based and learned evaluators.



\textbf{World Model Benchmarks.}
WorldModelBench~\cite{li2025worldmodelbench} and EWMBench~\cite{yue2025ewmbench} evaluate core dimensions like Perception, Prediction, and Generalization. 
Critically, as shown in Table~\ref{tab:final_comparison}, these benchmarks do not assess the crucial core dimension of Planning and Execution. Furthermore, no prior benchmark provides a metrics design that covers Planning Reasoning or Execution Accuracy in the robotics domain.
Our proposed \textbf{WoW-World-Eval} extends these efforts with a robotics-centric design, making it the most comprehensive embodied world model benchmark to date.

% . It is the first benchmark to provide a comprehensive evaluation across all five Core Dimensions. It is also the only benchmark to integrate these dimensions with a complete metric suite, including \texttt{VQ}, \texttt{IU}, \texttt{PL}, and the newly introduced \texttt{PR} and \texttt{EA} metrics